% ==========================================
% MỞ ĐẦU
% ==========================================
\chapter*{MỞ ĐẦU}
\addcontentsline{toc}{chapter}{MỞ ĐẦU}

\section*{1. Tính cấp thiết của đề tài}
Nhóm nghiên cứu có trách nhiệm trình bày đầy đủ các nội dung theo cấu trúc của báo cáo. Báo cáo hoàn chỉnh có độ dài từ 30 đến 40 trang A4, bao gồm hình ảnh, bảng biểu, số liệu và tài liệu tham khảo.

\section*{2. Phạm vi nghiên cứu}
Nội dung phạm vi nghiên cứu...

\section*{3. Phương pháp nghiên cứu}
\begin{itemize}
    \item Phương pháp khảo sát thực tiễn: Nếu có\ldots
    \item Phương pháp phân tích -- tổng hợp: Được sử dụng để xử lý số liệu\ldots
    \item Phương pháp mô phỏng hoặc thực nghiệm (nếu có): Mô phỏng quá trình\ldots
\end{itemize}

\section*{4. Cơ sở khoa học và ý nghĩa thực tiễn của đề tài}
\textbf{- Cơ sở khoa học:} Trình bày đầy đủ nội dung\ldots\\
\textbf{- Ý nghĩa thực tiễn:} Trình bày đầy đủ nội dung\ldots

\section*{5. Các nội dung nghiên cứu chính của đề tài}
Cấu trúc của báo cáo được chia thành 4 chương:
\begin{itemize}
    \item \textbf{Chương 1:} TỔNG QUAN VỀ ĐỀ TÀI
    \item \textbf{Chương 2:} CƠ SỞ LÝ THUYẾT VÀ CÔNG TRÌNH LIÊN QUAN
    \item \textbf{Chương 3:} PHƯƠNG PHÁP VÀ THIẾT KẾ HỆ THỐNG
    \item \textbf{Chương 4:} THỰC NGHIỆM VÀ ĐÁNH GIÁ
    \item \textbf{Chương 5:} KẾT LUẬN VÀ ĐỀ XUẤT HƯỚNG PHÁT TRIỂN
\end{itemize}% ==========================================
